%%%%%%%%%%%%%%%%%%%%%%%%%%%%%%%%%%%%%%%%%%%%%%%%%%%%%%%%%%%%%%%%%%%%%%%%%%%
%% UpdraftNet: Cooperative Urban Updraft Mapping for Energy-Aware
%%             UAV Delivery Swarms
%%
%% IEEE Conference Paper (6 pages, double-column)
%%
%% Text and layout follow the accepted PDF.  Everything that was changed
%% (review comments + numbers/setup aligned to the simulator) is typeset
%% in RED via \rd{...} or {\color{red}...}.  To get the final black
%% version, change \newcommand{\rd}[1]{{\color{red}#1}} to
%% \newcommand{\rd}[1]{#1} and replace \color{red} by nothing in displays
%% (or simply \renewcommand{\color}[1]{} -- see the line marked FINAL).
%%%%%%%%%%%%%%%%%%%%%%%%%%%%%%%%%%%%%%%%%%%%%%%%%%%%%%%%%%%%%%%%%%%%%%%%%%%

\documentclass[conference]{IEEEtran}
\IEEEoverridecommandlockouts

\usepackage[T1]{fontenc}
\usepackage[utf8]{inputenc}

\usepackage{amsmath,amssymb,amsfonts,amsthm}
\usepackage{algorithmic}
\usepackage{algorithm}
\usepackage{graphicx}
\usepackage{textcomp}
\usepackage{xcolor}
\usepackage{booktabs}
\usepackage{array}
\usepackage{mathtools}
\usepackage{cite}
\usepackage{placeins}

%% -- Marking of changes (red) ----------------------------------------
\newcommand{\rd}[1]{#1}   % red marking removed (final black version)
\newcommand{\rdisp}{}     % red marking removed (final black version)

%% -- Space optimisation (no content removed) ------------------------
\AtBeginDocument{%
  \setlength{\abovedisplayskip}{4pt plus 1pt minus 1pt}%
  \setlength{\belowdisplayskip}{4pt plus 1pt minus 1pt}%
  \setlength{\abovedisplayshortskip}{2pt plus 1pt}%
  \setlength{\belowdisplayshortskip}{2pt plus 1pt}%
}
\setlength{\abovecaptionskip}{3pt}
\setlength{\belowcaptionskip}{-3pt}
\setlength{\floatsep}{3pt plus 1pt minus 1pt}
\setlength{\textfloatsep}{5pt plus 1pt minus 1pt}
\setlength{\intextsep}{4pt plus 1pt minus 1pt}

%% -- Overflow guards ------------------------------------------------
\interdisplaylinepenalty=2500
\allowdisplaybreaks
\sloppy
\emergencystretch=1em

%% -- Shorthand that keeps wide displays inside the column -----------
\newcommand{\Tid}{T_{\mathrm{idle}}}
\newcommand{\wdet}{w_{\mathrm{det}}}

\newtheorem{theorem}{Theorem}
\newtheorem{lemma}{Lemma}
\newtheorem{proposition}{Proposition}
\newtheorem{corollary}{Corollary}
\newtheorem{property}{Property}
\newtheorem{definition}{Definition}

\begin{document}

\title{UpdraftNet: Cooperative Urban Updraft Mapping\\
       for Energy-Aware UAV Delivery Swarms}

%% The accepted PDF shows no author block; restore the lines below if needed.
\author{\IEEEauthorblockN{Karan Yadav, Adarsh Narayan, Prajjwal Patel,
        and Arijit Roy,~\IEEEmembership{Senior Member,~IEEE}}
\IEEEauthorblockA{Department of Computer Science and Engineering\\
Indian Institute of Technology Patna, Patna 801106, India\\
Email: \{karan\_2511ai24, arijitroy\}@iitp.ac.in, 24155637@kiit.ac.in, adarshnarayan517@gmail.com}
}

\maketitle

%% ------------------------------------------------------------------
\begin{abstract}
In this work, we propose UpdraftNet, a cooperative framework that enables
a UAV swarm to jointly discover, share, and exploit urban updraft structure
for energy-aware routing. Urban delivery UAVs fly constant-throttle
trajectories past sizable buildings whose windward faces produce 1--3\,m/s
of upward vertical flow~\cite{holmes,asce,kawamura}, yet no current delivery
system exploits this freely available lift. UpdraftNet integrates
an onboard vertical-wind estimator that fuses barometric altitude
rate, vertical inertial acceleration, and rotor-torque signals. The
proposed system also includes a Bayesian spatiotemporal map
with explicit variance decay, a constant-size \rd{one-hop peer-exchange (gossip-style) protocol whose payload is
independent of swarm size}, and an updraft-biased A$^\star$ planner that minimises cost function routing cost
over the \rd{UAV's fused} map. We establish that coverage time scales
as $O(1/n)$ with swarm size, the planned trajectory satisfies
a routing-optimality bound relative to true energy cost, the
variance-discounted wind estimate decays toward zero as idle
time grows, and a detour \rd{lowers the routing cost} if and only if the
updraft gain exceeds the extra path length. Simulations over 60
urban configurations with up to 16 UAVs show that UpdraftNet
achieves \rd{16.1\% energy saving at $n=16$ over shortest-path routing under the stated power model, about $1.8\times$ that of UAVs that do not share their maps,} and reduces \rd{fleet-union} map coverage time by up to \rd{94\%} compared to a
single UAV operating alone.
\end{abstract}
\begin{IEEEkeywords}
UAV swarm, energy-aware routing, urban air
mobility, cooperative mapping, dynamic soaring, updraft exploitation.
\end{IEEEkeywords}

%% ------------------------------------------------------------------
\section{Introduction}

Multirotor UAVs carry sufficient battery for 20--30 minutes
of flight~\cite{shakhatreh}. UAV delivery missions consist of distinct flight
phases, takeoff, level (cruise) flight, hovering, and landing,
and the total energy demand is determined by the power draw
and duration of each phase, with the cruise segment typically
dominating total mission energy~\cite{kirschstein}. Reducing cruise energy
increases deliveries per charge and slows battery degradation.
sizable buildings deflect wind upward on their windward face
and downward on their leeward side~\cite{holmes,asce}. This vertical
flow reaches 1--3\,m/s depending on building height and wind
speed~\cite{holmes,asce,kawamura}. Birds exploit similar flow near cliffs and
coastal buildings to stay aloft with less energy~\cite{pooley}. Commercial
multirotor UAVs fly at constant throttle and ignore this flow
entirely.

Two reasons explain the gap between available updraft
energy and current UAV practice. First, a single UAV covers
only a small portion of the urban airspace at any time~\cite{la},
and building-induced updrafts shift on timescales of tens of
seconds~\cite{tominaga}, therefore old readings lose value fast. Second,
existing work focuses on rural convective thermals~\cite{allen,reichmann}.
Rural thermals differ in scale and structure from building-induced
updrafts and do not require multi-UAV coordination.
We propose UpdraftNet, where each UAV estimates local
vertical wind from its existing onboard sensors and \rd{exchanges its measurements with a randomly chosen peer through a lightweight message}. \rd{Each UAV fuses them into its own Bayesian wind
map.} Each UAV then plans routes that favour updraft zones
and avoid downdraft zones.

\subsection{Contributions}
This work makes the following \textit{contributions}:
\begin{itemize}
\item We formulate cooperative urban updraft exploitation as
a joint estimation and routing problem and introduce
a Bayesian spatiotemporal updraft map with explicit
confidence decay.
\item We design a constant-size \rd{one-hop peer-exchange protocol} and an updraft-biased route
planner for multi-UAV 3D mapping. Both fit
small-UAV memory and latency budgets.
\item We prove four analytical results, a coverage-time bound,
a routing-optimality bound, a freshness-decay characterisation,
and a profitable-detour condition, showing \rd{that map sharing empirically raises individual energy saving by about $80\%$}. Simulation over 60 urban configurations with up to
16 UAVs \rd{reports 16.1\%} energy saving and coverage
time falling from \rd{$603\,\text{s}$} at $n=1$ to \rd{$37\,\text{s}$} at $n=16$.
\end{itemize}

%% ------------------------------------------------------------------
\section{Related Work}

This section discusses existing work related to our proposed
work. We divide it into three parts, UAV soaring and energy
harvesting, urban UAV path planning, and cooperative UAV
field mapping.

\subsection{UAV Soaring and Energy Harvesting}

Due to the abundance of thermal and flow energy in the
atmosphere, UAVs are studied as platforms for autonomous
aerodynamic energy harvesting. Allen~\cite{allen} demonstrated autonomous thermal soaring over rural terrain, establishing the
first proof of concept that an unmanned fixed-wing aircraft can
sustain flight by riding convective thermals without propulsion. Along the same direction, Pooley~\textit{et~al.}~\cite{pooley} analysed
multi-UAV settings in which agents share thermal-location
indicators using hawk-inspired flocking behaviour, amplifying
individual lift harvesting through inter-agent communication.
More broadly, Wei~\textit{et~al.}~\cite{wei} surveyed solar and mechanical
energy-harvesting mechanisms for UAVs, including dynamic
soaring, wind-induced vibration, and flapping-wing motion,
and discussed integration challenges for delivery platforms. In
these works, however, the authors addressed rural thermals and
did not consider building-induced urban updrafts, which differ
in spatial scale, directional structure, and the coordination
demands they impose on a delivery swarm operating under
a mission deadline.

\subsection{Urban UAV Path Planning}

Urban UAV route planning studies incorporated secondary
objectives such as noise minimisation and ground-risk reduction. Sarhan~\textit{et~al.}~\cite{sarhan} proposed a reinforcement-learning path
planner that minimises noise exposure in urban corridors.
Energy appeared implicitly through path length and vertical
wind did not enter the cost function. Similarly, Chen~\textit{et~al.}~\cite{chen}
incorporated ground-risk and noise-mitigation terms into a
joint flight-path optimisation for urban air mobility. Along
the same line, Tan~\textit{et~al.}~\cite{tan} designed a low-noise trajectory
planner using a noise assessment simulator. On the other hand,
Aldao~\textit{et~al.}~\cite{aldao} addressed obstacle avoidance in dynamic
building environments using a point-cloud-derived 3-D geometric algorithm instead of a noise or risk objective. None of
these works modelled the aerodynamic structure of the vertical
wind field. Updraft exploitation was therefore not achievable
in any of them.

\subsection{Cooperative UAV Field Mapping}

UAV swarms estimated distributed scalar fields such as
pollutant concentration, radiation, and RF signal strength. La
and Sheng~\cite{la} developed a distributed sensor-fusion approach for
scalar-field mapping using mobile sensor networks. However, the shared field was passive and did not affect the
energetics of the sensing UAVs. In one such application,
Bayerlein~\textit{et~al.}~\cite{bayerlein} proposed multi-UAV path planning for
wireless data harvesting using deep reinforcement learning.
The shared quantity there was channel state, not a physical
wind field. On a more theoretical level, Boyd~\textit{et~al.}~\cite{boyd} established the
convergence properties of randomised gossip algorithms. These properties \rd{motivate} the
\rd{peer-exchange} protocol used in this work.

\textit{\textbf{Synthesis.}}
No prior work jointly mapped building-induced urban updrafts and
used that map to reduce routing energy for a delivery swarm.
UpdraftNet fills this gap.

%% ------------------------------------------------------------------
\section{System and Field Model}
\label{sec:model}

We consider a set $\mathcal{U}=\{U_1,\ldots,U_n\}$ of $n$ multirotor
UAVs operating in a bounded urban region $\mathcal{R}\subseteq\mathbb{R}^3$
over a finite mission horizon $[0,T]$, where $T$ denotes the maximum mission
duration. Each UAV executes one or more delivery missions,
each defined by an origin--destination pair and a deadline. The
ambient wind field, after deflection by buildings, has a vertical
component $w(\mathbf{x},t)$ at position $\mathbf{x}\in\mathcal{R}$ and time $t\in[0,T]$. By
convention, $w>0$ denotes upward flow.

Let $\{B_1,\ldots,B_M\}$ denotes the set of buildings in $\mathcal{R}$. For
each building $B_m$ with height $H_m$ and upwind face normal
$\hat{\mathbf{n}}_m$, ambient wind $\mathbf{u}_\infty(t)$ defines a
windward updraft cone $\Omega_m^+$ and a leeward downdraft cone $\Omega_m^-$. Inside $\Omega_m^+$, building-engineering practice~\cite{holmes,asce} models the mean vertical wind as
$w_m^+(\mathbf{x},t)$ as:
\begin{equation}
w_m^+(\mathbf{x},t) =
\alpha\,\langle \mathbf{u}_\infty(t),\hat{\mathbf{n}}_m\rangle\,
\phi\!\left(\frac{H_m}{H_{\mathrm{ref}}}\right)\,
\psi(d(\mathbf{x},B_m),z)
\label{eq:windward}
\end{equation}
where $\alpha \in (0,1)$ represents an empirical deflection coefficient
derived from published wind-tunnel data~\cite{kawamura,tominaga},
$\phi(\cdot)$ is a height-scaling
function. Here, $H_{\mathrm{ref}}$, $d(\mathbf{x},B_m)$, and $\psi(d,z)\in[0,1]$ denote a baseline reference height, the perpendicular distance from $\mathbf{x}$ to
the windward face, and  a spatial-decay envelope
that falls to zero beyond a distance proportional to $H_m$ respectively. The leeward
field $w_m^-$ follows an analogous expression, substituting $\alpha$ with a negative leeward deflection coefficient $-\alpha_d$.

\begin{figure}[!t]
\centering
\includegraphics[width=0.66\columnwidth]{fig_geometry.png}
\caption{UAV updraft geometry and UpdraftNet pipeline.}
\label{fig:geometry}
\end{figure}

The total vertical field is the superposition $w(\mathbf{x}, t)$ is:
\begin{equation}
w(\mathbf{x},t) = \sum_{m=1}^{M}\!\Big[
w_m^+(\mathbf{x},t)\,\mathbf{1}_{\mathbf{x}\in\Omega_m^+}
+w_m^-(\mathbf{x},t)\,\mathbf{1}_{\mathbf{x}\in\Omega_m^-}\Big]
+\eta(\mathbf{x},t)
\label{eq:field}
\end{equation}
where $\eta(\mathbf{x},t)$ collects turbulent fluctuations. We treat $\eta$
as a zero-mean process with locally bounded variance $\sigma_\eta^2(\mathbf{x})$
and finite spatial correlation length $\ell$,
consistent with published urban canyon flow measurements~\cite{britter}. \rd{Let $\wdet=w-\eta$ be the bracketed sum in Eq.~\eqref{eq:field}; the map estimates $\wdet$ and the energy model below uses it.}

A UAV at position $\mathbf{x}_i(t)$ does not observe
$w(\mathbf{x}_i(t),t)$ directly but estimates it from three existing onboard
signals, barometric altitude rate $\dot{h}_i^{\text{baro}}(t)$, vertical
inertial acceleration $a_{iz}(t)$, and rotor torque $\tau_i(t)$ via a
thrust model $f_z(\tau_i)$. Newton's second law applied to the vertical axis gives the following expression for conservation of vertical momentum
on a multirotor of mass $m_i$ such that:
\begin{equation}
m_i\bigl(a_{iz}+g\bigr) = f_z(\tau_i) + C_z\,w(\mathbf{x}_i,t)
\label{eq:vert_balance}
\end{equation}
where $g$ is gravitational acceleration and $C_z>0$ is the vertical
aerodynamic coupling coefficient. We represent Eq.~\eqref{eq:vert_balance} as:
\begin{equation}
\hat{w}_i(t) = \frac{1}{C_z} \Big[ m_i \bigl(a_{iz} + g \bigr) - f_z \bigl(\tau_i(t) \bigr) \Big]
\label{eq:w_estimator}
\end{equation}
We assume an additive Gaussian model $\hat{w}_i=w(\mathbf{x}_i,t)+\xi_i$ with
$\xi_i\sim\mathcal{N}(0,\sigma_w^2)$, where the noise variance $\sigma_w^2$ is characterised
via offline sensor calibration routines~\cite{thrun}. \rd{Since $\hat{w}_i=\wdet+\eta_i+\xi_i$ with $\eta_i,\xi_i$ independent, one sample of $\wdet$ has variance $R=\sigma_w^2+\sigma_\eta^2$.}

We discretise $\mathcal{R}$ into voxels of side \rd{$\Delta$} meters. For voxel $k$
at update step $\tau$, $\hat{w}_k^{(\tau)}$ and $\sigma_k^{2,(\tau)}$ denote the running estimate
and its variance, respectively. When a measurement $\hat{w}_i$ falls
within voxel $k$, the standard Bayesian recursion~\cite{lawrance} updates
the map using the following equations:
\begin{align}
\hat{w}_k^{(\tau)} &=
\frac{\rdisp R\,\hat{w}_k^{(\tau-1)}
      + \sigma_k^{2,(\tau-1)}\,\hat{w}_i}
     {\rdisp R+\sigma_k^{2,(\tau-1)}}
\label{eq:bayes_mean}\\
\sigma_k^{2,(\tau)} &=
\frac{\rdisp R\,\sigma_k^{2,(\tau-1)}}
     {\rdisp R+\sigma_k^{2,(\tau-1)}}
\label{eq:bayes_var}
\end{align}
In the absence of fresh sensor data, variance grows because the
underlying field is non-stationary. We implement a first-order
inflation step defined in Eq.~\eqref{eq:inflation} to capture the increase in map uncertainty,
\begin{equation}
\sigma_k^{2,(\tau)} \leftarrow \sigma_k^{2,(\tau)}
   + \rho\,\delta_t\,\sigma_\eta^2
\label{eq:inflation}
\end{equation}
where \rd{$\delta_t$ is the time since the previous update and} $\rho>0$ \rd{(in s$^{-1}$)} is the temporal-volatility rate of the wind field.
Eqs.~\eqref{eq:bayes_var} and \eqref{eq:inflation} together produce the \rd{monotone} confidence
decay analysed in Theorem~\ref{thm:decay}.

%% ------------------------------------------------------------------
\section{The UpdraftNet Framework}
\label{sec:framework}

\subsection{Peer Broadcast Protocol}

Every $\delta_b$ seconds, each UAV selects a random peer and
transmits a state update message. The message contains the
voxel index, the \rd{UAV's latest raw measurement $\hat{w}_i$, and its sample variance $R$}
for the voxel $k$ visited during the previous interval. Since
the message is a fixed-size triple \rd{$(k,\hat{w}_i,R)$} containing a
single voxel index and two scalars, per-broadcast payload size
is constant and independent of swarm size $n$ or map resolution
$K$. Upon receiving the data, peer UAV $j$ treats the incoming
\rd{measurement} as a fresh measurement and fuses the value into the
local map using Eqs.~\eqref{eq:bayes_mean} and \eqref{eq:bayes_var}.
\rd{The exchange is one-hop (received measurements are not relayed, so none is counted twice), does not
reproduce a centralised map, and convergence to a common map is not claimed;
the scheme keeps} $O(1)$ per-node communication overhead and $O(|\mathcal{R}|/\rd{\Delta}^3)$
memory.

\subsection{Updraft-Biased Cost Function}

Let $\Gamma=(\mathbf{x}(s))_{s\in[0,L]}$ be a candidate flight trajectory of arc
length $L$. The cumulative cruise energy of a multirotor along
$\Gamma$ is modelled as:
\begin{equation}
E(\Gamma)=\rdisp\frac{P_0}{V}\int_0^L\!\Bigl[1-\min\{\beta\,\wdet(\mathbf{x}(s),t),\,0.9\}\Bigr]ds
\label{eq:energy}
\end{equation}
where $P_0$ is the parasitic-plus-induced power at level cruise
\rd{and} $V$ is the cruise airspeed\rd{; $\beta$ [s/m] is a model parameter: rotor-disk theory gives $P\approx P_0(1-\beta w)$ with
$\beta\le m_ig/P_0$~\cite{leishman} ($\approx0.08\,\mathrm{s/m}$ for $m_i=1\,\mathrm{kg}$, $P_0=120\,\mathrm{W}$), and we use $\beta=1.5\,\mathrm{s/m}$, deliberately optimistic.
The product $\beta \wdet$ is capped at $0.9$ ($P\ge0.1P_0$), downdrafts are not capped, and the planner and the energy evaluation share $\beta$ and the cap, so the savings
below are savings under this power model}.
We \rd{replace $\wdet$ in} Eq.~\eqref{eq:energy} \rd{by the estimate $\bar{w}$} to obtain the cost
function:
\begin{equation}
J(\Gamma)=\int_0^L\!\bigl[1 - \rdisp\min\{\beta\,\bar{w}(\mathbf{x}(s),\hat{w},\sigma_k^2),\,0.9\}\bigr]\,ds
\label{eq:cost}
\end{equation}
where $\bar{w}$ is a
variance-discounted point estimate,
\begin{equation}
\bar{w}(\mathbf{x},\hat{w},\sigma^2)=
\hat{w}\cdot\!\left(1 - e^{-\nu/\sigma^2}\right)
\label{eq:variance_discount}
\end{equation}
such that $\nu>0$\rd{ has the units of $\sigma^2$ ($\mathrm{m^2/s^2}$)}. Eq.~\eqref{eq:variance_discount} shrinks the effective
estimate when the map is uncertain and approaches the
raw estimate when the map is confident. \rd{Positive $\wdet$ averages 0.65\,m/s (0.84\,m/s on UpdraftNet routes), peaking at 2.28\,m/s, and $\beta \wdet>0.9$ for $\wdet>0.6\,\mathrm{m/s}$ (45\% of updraft voxels, 62\% on routes), so the cap is active.}

\begin{algorithm}[h]
\caption{UpdraftNet (per UAV $U_i$, per scheduling tick)}
\label{alg:updraftnet}
\footnotesize
\begin{algorithmic}[1]
\REQUIRE Current position $\mathbf{x}_i$, destination $\mathbf{x}_g$,
local map $\mathcal{M}_i$, broadcast bus $\mathcal{B}$.
\STATE Read onboard signals $(\dot h_i^{\text{baro}}, a_{iz}, \tau_i)$.
\STATE Compute $\hat{w}_i$ from Eq.~\eqref{eq:w_estimator}.
\STATE Update voxel $k(\mathbf{x}_i)\in\mathcal{M}_i$ via
Eqs.~\eqref{eq:bayes_mean} and \eqref{eq:bayes_var}.
\STATE \rd{Send} triple \rd{$(k,\hat{w}_i,R)$ to one random peer via} $\mathcal{B}$.
\STATE Receive peer triples. For each triple, repeat step~3.
\STATE Apply temporal inflation Eq.~\eqref{eq:inflation} to all voxels.
\IF{re-plan timer expired}
  \STATE Compute $\Gamma^\star\!=\!\arg\min_\Gamma J(\Gamma)$ on $G$
  via A$^\star$ over cost Eq.~\eqref{eq:cost}.
  \IF{$J(\Gamma^\star)$ saves $\geq \theta_E$ vs.\ current plan}
    \STATE Commit to $\Gamma^\star$ subject to flight-envelope checks.
  \ENDIF
\ENDIF
\RETURN Executed control along committed plan.
\end{algorithmic}
\end{algorithm}

\subsection{Routing}

We discretise the navigable airspace into a graph $G=(\mathcal{V},\mathcal{E})$
whose vertices are voxel centres and whose edges
connect voxels within a kinematic-feasibility radius \rd{(26-neighbour edges)}. Each edge
receives a cost equal to its arc-length contribution to Eq.~\eqref{eq:cost}\rd{, evaluated at the mean of its end-point estimates}.
The optimal path is found by A$^\star$ search~\cite{hart} with a Euclidean
lower-bound heuristic \rd{$h(\mathbf{x})=(1-\min\{\beta\bar{w}_{\max},0.9\})\|\mathbf{x}-\mathbf{x}_g\|\ge0$, with $\bar{w}_{\max}$ the largest non-negative map estimate,}
that is admissible under Eq.~\eqref{eq:cost} \rd{since every edge costs at least $1-\min\{\beta\bar{w}_{\max},0.9\}\ge0.1$ per metre. The multi-UAV simulations use Dijkstra's algorithm on $G$ (same optimum); A$^\star$ is used in the latency tests.}

%% ------------------------------------------------------------------
\section{Analysis}
\label{sec:analysis}
We analyse the proposed scheme, UpdraftNet, using coverage time, routing
optimality, freshness decay and energy-detour tradeoff, defined below.

\subsection{Coverage Time}

Let $K=|\mathcal{R}|/\rd{\Delta}^3$ be the total voxel count and $\theta>0$ be
a target confidence threshold. The value of $\theta$ is provided in
simulation results.

\begin{property}
\label{prop:coverage}
If each UAV visits voxels i.i.d.\ uniformly over $\mathcal{R}$ and
\rd{measures every $\delta_m$ seconds}, the time \rd{for the fleet (the union of the UAVs' local maps)} to bring $K$ voxels
below variance threshold $\theta$ \rd{(ignoring the inflation of Eq.~\eqref{eq:inflation})} scales as
\begin{equation}
T_{\theta}(n) \rd{\simeq} \frac{K\,\rdisp\delta_m}{n}\,\rdisp
\bigl(\ln K+(N_\theta-1)\ln\ln K\bigr),
\label{eq:coverage}
\end{equation}
by a coupon-collector argument~\cite{motwani,doerr}\rd{ (double Dixie cup~\cite{newman}),}
reducing $\sigma_k^2$ below $\theta$ needs \rd{$N_\theta=\lceil R(1/\theta-1/\sigma_0^2)\rceil$ measurements per voxel ($\sigma_0^2$ the prior variance)}, and $n$ UAVs
jointly generate \rd{$nT/\delta_m$} measurements over $K$ voxels, giving
the $O(1/n)$ leading term and \rd{the $\ln K+(N_\theta-1)\ln\ln K$ factor} from
maximum-load fluctuation.
\end{property}

\rd{Eq.~\eqref{eq:coverage} bounds the fleet-union map under i.i.d.\ uniform visits, whereas the simulations use goal-directed routes, so only the $1/n$ trend is compared.
A local map ingests $\approx1/\delta_m+1/\delta_b=1.5$ measurements/s for any $n$, so local-map coverage does not scale as $1/n$.}

\subsection{Routing Optimality}

\begin{property}[Routing Suboptimality]
\label{prop:routing}
\rd{Fix the wind field at the planning instant $t$. Let $\varepsilon=\sup_{\mathbf{x}}|\wdet(\mathbf{x},t)-\bar{w}(\mathbf{x})|$ and let $L_{\max}$ bound the arc length of candidate paths in $G$.}
The trajectory $\Gamma^\star$ returned by Algorithm~\ref{alg:updraftnet}
satisfies
\begin{equation}
E(\Gamma^\star) \leq \min_{\Gamma\rd{\in G}}\,E(\Gamma) \rd{{}+ \frac{2P_0\,\beta\,\varepsilon\,L_{\max}}{V}}.
\label{eq:routing_opt}
\end{equation}
This follows because the planner does not optimise the true
energy $E$ directly. It optimises a first-order surrogate $J$ built
around the estimated wind field $\bar{w}$\rd{. With $J_w$ the surrogate evaluated at the true wind, $E(\Gamma)=\frac{P_0}{V}J_w(\Gamma)$, and since
$u\mapsto\min\{u,0.9\}$ is $1$-Lipschitz, $|J_w(\Gamma)-J(\Gamma)|\le\beta\varepsilon L_{\max}$ for every $\Gamma\in G$}.
Since A$^\star$ solves the surrogate problem exactly, it
cannot do worse on $J$ than the true minimum-energy trajectory
$\Gamma^{\mathrm{opt}}$ does, so $J(\Gamma^\star)\le J(\Gamma^{\mathrm{opt}})$.
\rd{Hence $J_w(\Gamma^\star)\le J(\Gamma^\star)+\beta\varepsilon L_{\max}\le J(\Gamma^{\mathrm{opt}})+\beta\varepsilon L_{\max}\le J_w(\Gamma^{\mathrm{opt}})+2\beta\varepsilon L_{\max}$,
and multiplying by $P_0/V$ gives Eq.~\eqref{eq:routing_opt}.} The
bound holds \rd{under the linear level-cruise model, i.e.,} when
higher-order aerodynamic terms stay negligible.
\end{property}

\rd{The bound is worst-case ($\varepsilon$ is not small by construction, as unvisited voxels have $\bar{w}=0$): at $\beta=1.5\,\mathrm{s/m}$ each $0.1\,\mathrm{m/s}$ of $\varepsilon$ adds
$2\beta\varepsilon=30\%$ of $P_0L/V$, so the benefit is shown by simulation.} The bound tightens as
map quality improves\rd{; $\varepsilon$ is the error of the planning UAV's local map, set by its ingestion rate rather than $n$, consistent with the saving saturating from $n=8$ to $16$}.

\subsection{Freshness Decay}

\begin{theorem}
\label{thm:decay}
Suppose voxel $k$ receives no measurement over an
interval of length $\Tid$, during which the wind has temporal
volatility rate $\rho$. Then
\begin{equation}
\sigma_k^{2}(\tau+\Tid) =
\sigma_k^{2}(\tau) + \rho\,\Tid\,\sigma_\eta^2
\label{eq:decay_linear}
\end{equation}
and the variance-discounted estimate satisfies
\begin{equation}
\frac{\bigl|\bar{w}(\Tid)\bigr|}{|\hat{w}|}
= 1 - \exp\!\left(
\frac{-\nu}{\sigma_k^{2}(\tau)+\rho\,\Tid\,\sigma_\eta^2}\right)
\label{eq:decay_envelope}
\end{equation}
which decays toward zero as $\Tid\to\infty$.
\end{theorem}

\begin{IEEEproof}
Equation Eq.~\eqref{eq:decay_linear} follows from Eq.~\eqref{eq:inflation} by
induction. Substituting into Eq~\eqref{eq:variance_discount}
yields Eq.~\eqref{eq:decay_envelope}. The exponent's
denominator grows linearly with $\Tid$, so the discounted estimate falls toward zero.
\end{IEEEproof}

Theorem~\ref{thm:decay} gives a
principled trigger for re-measurement. With \rd{$\sigma_k^2(\tau)=\theta\approx0.0506\,\mathrm{m^2/s^2}$,} $\sigma_\eta^2=0.04\,\mathrm{m^2/s^2}$
and $\rho=0.02$\rd{$\,\mathrm{s^{-1}}$}, the discount \rd{ratio starts at $0.63$ and} falls to \rd{$0.5$} within approximately
\rd{$27\,\mathrm{s}$}.

\subsection{Energy--Detour Tradeoff}

\begin{theorem}
\label{thm:detour}
Let $\Gamma_\mathrm{sp}$ be the shortest path of length $L_\mathrm{sp}$, and
let $\Gamma_\mathrm{up}$ be a candidate path of length $L_\mathrm{up}=L_\mathrm{sp}+\Delta L$ that
passes through an updraft zone of mean upward flow $\bar{w}^+$ over
a sub-arc of length $L^+$\rd{. Let $\bar g=\frac{1}{L^+}\int_{\text{sub-arc}}\min\{\beta\bar{w}(\mathbf{x}(s)),0.9\}\,ds$ (so $\bar g\le\min\{\beta\bar{w}^+,0.9\}$, with equality for constant $\bar w$) and assume $\bar{w}=0$ on $\Gamma_\mathrm{sp}$ and on
$\Gamma_\mathrm{up}$ outside that sub-arc}. The detour \rd{lowers the routing cost $J$} if and
only if
\begin{equation}
\rdisp\bar g\,L^+ > \Delta L
\label{eq:detour_condition}
\end{equation}
or equivalently, $\Delta L/L^+ < \rd{\bar g}$\rd{; since $\bar g\le0.9$, any profitable detour needs $\Delta L<0.9L^+$}.
\end{theorem}

\begin{IEEEproof}
Substituting $\Gamma_\mathrm{up}$ into Eq.~\eqref{eq:cost} and subtracting
$J(\Gamma_\mathrm{sp})$ gives
\rd{$J(\Gamma_\mathrm{up})=(L_\mathrm{up}-L^+)+(1-\bar g)L^+=L_\mathrm{sp}+\Delta L-\bar g L^+$, so} $J(\Gamma_\mathrm{up})-J(\Gamma_\mathrm{sp})=\Delta L - \rd{\bar g}\,L^+$, the detour
reduces $J$ iff this difference is negative, which is Eq.~\eqref{eq:detour_condition}.
\end{IEEEproof}

\begin{figure*}[!t]
\centering
\includegraphics[width=0.65\textwidth]{fig_simulation_v5.png}
\caption{Simulation results for UpdraftNet and baseline schemes.
(a)~Coverage time to map 40\% of voxels vs.\ swarm size\rd{ (dashed: Poisson reference)}.
(b)~\rd{Simulated} energy saving vs.\ ambient wind speed \rd{($n=8$)}.
(c)~\rd{Simulated} re-planning latency vs.\ grid size.
(d)~Per-node \rd{communication} overhead vs.\ swarm size \rd{(all-to-all: same message to every peer)}.
(e)~Wind-map confidence decay over idle time \rd{(line: two measurements; dots: Monte-Carlo, two to four; MC points start at $0.75$ vs.\ $0.63$ because they fuse two to four measurements)}.
(f)~\rd{Simulated} energy saving vs.\ delivery tasks \rd{per UAV} ($n=8$).}
\label{fig:results}
\end{figure*}

%% ------------------------------------------------------------------
\section{Performance Analysis}
\label{sec:evaluation}

\subsection{Simulator Setup}
This section discusses the performance of the proposed
scheme, UpdraftNet. We model the simulated airspace as a
$12\!\times\!12\!\times\!3$ voxel grid ($\Delta=20\,\text{m}$, $K=432$ voxels) representing
a $240\,\text{m}\!\times\!240\,\text{m}\!\times\!60\,\text{m}$ urban volume. Eight buildings \rd{with $20$--$40\,\text{m}$ footprints and} of height
$\in[40,100]\,\text{m}$ are placed randomly per trial. \rd{In Eq.~\eqref{eq:windward}, $\phi(r)=\min(r,2)$, $\psi=e^{-d/0.6H_m}$ ($d<H_m$), leeward $e^{-d/0.8H_m}$ ($d<1.5H_m$), with a Gaussian lateral decay.} Each UAV cruises
at $V=10\,\text{m/s}$\rd{, samples at 1\,Hz ($\delta_m=1\,\text{s}$), flies random deliveries of at least $80\,\text{m}$, re-plans every $10\,\text{s}$ ($\theta_E=2\%$),} and \rd{sends its newest raw measurement to one random peer (one hop)} every $\delta_b=2\,\text{s}$.
\rd{Each sample is $\hat{w}=\wdet+\eta+\xi$ with $\eta\sim\mathcal{N}(0,\sigma_\eta^2)$ drawn per sample, so its variance is $R=\sigma_w^2+\sigma_\eta^2=0.1025\,\text{m}^2/\text{s}^2$ (the estimator of Eqs.~\eqref{eq:vert_balance}--\eqref{eq:w_estimator} is abstracted by its output noise).}
A voxel is
classified as confidently mapped \rd{once at least one UAV's map holds} at least two \rd{measurements of it, regardless of age (fleet-union coverage; Eq.~\eqref{eq:inflation} is ignored)}.
\rd{With the prior $\sigma_0^2=4\,\text{m}^2/\text{s}^2$ we define the confidence threshold $\theta\equiv1/(1/\sigma_0^2+2/R)$ ($=0.0506\,\text{m}^2/\text{s}^2$ after rounding), the posterior variance after two samples, so that $N_\theta=2$ exactly}. Coverage is declared when $40\%$ of \rd{navigable} voxels satisfy
this criterion\rd{ within $900\,\text{s}$}.
\rd{Energy is evaluated with the deterministic wind $\wdet$ (without $\eta$) via $P=P_0(1-\min\{\beta \wdet,0.9\})$, as in the planner, and compared with the shortest path. Single-drone soaring is the planner without sharing; noise-aware routing weights length by population-based noise exposure.}

\begin{table}[!t]
\renewcommand{\arraystretch}{1.0}
\caption{Simulation Parameters}
\label{tab:params}
\centering
\begin{tabular}{ll}
\toprule
\textbf{Parameter} & \textbf{Value} \\
\midrule
Voxel size $\Delta$             & $20\,\text{m}$                                  \\
Grid dimensions                  & $12\times12\times3$                             \\
Ambient wind $u_\infty$          & $7\,\text{m/s}$~\cite{asce}                     \\
Deflection coeff.\ $\alpha$      & $0.15$~\cite{kawamura}                          \\
Leeward coeff.\ $\alpha_d$       & $0.35$~\cite{holmes}                            \\
Reference height $H_\text{ref}$  & $50\,\text{m}$~\cite{holmes}                    \\
Estimator noise $\sigma_w$       & $0.25\,\text{m/s}$ (IMU datasheet)              \\
Turbulence var.\ $\sigma_\eta^2$ & $0.04\,\text{m}^2/\text{s}^2$ (assumed)         \\
Temporal volatility $\rho$       & $0.02$\rd{$\,\text{s}^{-1}$}                    \\
Cruise speed $V$                 & $10\,\text{m/s}$                                \\
Cruise power $P_0$               & $120\,\text{W}$                                 \\
Routing weight $\beta$           & $1.5$\rd{$\,\text{s/m}$ (empirical)} \\
Broadcast interval $\delta_b$    & $2\,\text{s}$                                   \\
\rd{Discount sharpness $\nu$; prior $\sigma_0^2$} & \rd{$0.05$; $4\,\text{m}^2/\text{s}^2$} \\
Trials per config.\rd{; deliveries per UAV} & $60$\rd{; $24$}                                  \\
\bottomrule
\end{tabular}
\end{table}

\subsection{Results and Discussion}

Fig.~\ref{fig:results}(a) compares the coverage time achieved by different
swarm sizes. As the swarm size increases from 1 to 16, the
coverage time decreases from \rd{$603\,\text{s}$} to \rd{$37\,\text{s}$}, \rd{consistent with the $1/n$ trend of} Property~\ref{prop:coverage}.
\rd{The dashed line is the Poisson estimate $\lambda^\star K_{\mathrm{nav}}\delta_m/n$ ($K_{\mathrm{nav}}\approx385$, $1-e^{-\lambda^\star}(1+\lambda^\star)=0.4$, $\lambda^\star=1.376$), not Eq.~\eqref{eq:coverage};
single-drone soaring covers as fast because coverage is fleet-level.}

Fig.~\ref{fig:results}(b) compares UpdraftNet with shortest-path routing,
noise-aware routing, and single-UAV soaring under different
ambient wind speeds. UpdraftNet consistently achieves the
highest energy savings across all wind conditions\rd{, rising from $5.6\%$ at $3\,\text{m/s}$ to $22.0\%$ at $13\,\text{m/s}$,} because
stronger ambient wind raises the updraft estimate at affected
voxels and lowers their routing cost, an effect neither baseline
can exploit, shortest-path routing since wind never enters its
cost function, and single-UAV soaring since its map covers
only a fraction of the available updraft zones.

Fig.~\ref{fig:results}(c) plots route re-planning latency against grid size.
UpdraftNet \rd{needs about} $5\,\text{ms}$ at $K=432$ voxels \rd{(Python A$^\star$ on one $2.8\,\text{GHz}$ core)} because its
planner reads one precomputed wind value per voxel instead
of evaluating a cost model along each edge. \rd{Noise-aware and shortest-path routing are faster ($1.8$ and $0.3\,\text{ms}$ at $K=400$) because their heuristics are tighter; all grow roughly linearly with grid size.}

Fig.~\ref{fig:results}(d) shows per-node broadcast overhead across swarm
sizes. Overhead stays flat at \rd{$5\,\text{bytes/s}$ (one 10-byte message per $\delta_b$)} for all \rd{$n\ge2$}, confirming
the $O(1)$ \rd{one-hop exchange} property, since each \rd{message} carries a fixed
voxel index and two scalars regardless of swarm size. \rd{An all-to-all broadcast of the same message}
\rd{reaches $75\,\text{bytes/s}$ at $n=16$ because its traffic grows linearly with $n$}, making it unsuitable for large fleets.

\begin{table}[!t]
\renewcommand{\arraystretch}{1.0}
\caption{Comparison with Baselines ($n=8$/$n=16$, 60 trials).
Map coverage rate is the fraction of trials achieving 40\% voxel coverage
within $900\,\text{s}$\rd{ (others count as $900\,\text{s}$ in the mean)}.}
\label{tab:comparison}
\centering
\begin{tabular}{lccc}
\toprule
\textbf{Scheme} & \textbf{Energy} & \textbf{Cov.\ time} & \textbf{Map cov.} \\
                & \textbf{saving} &                     & \textbf{rate}     \\
\midrule
Shortest path (baseline)        & $0.0\%$           & N/A             & $0\%$   \\
Single soaring ($n=1$)          & \rd{$9.8\%$}           & \rd{$603\,\text{s}$} & \rd{$98\%$}  \\
\rdisp Single soaring ($n=8$)     & \rd{$9.2\%$}           & \rd{$73\,\text{s}$}  & \rd{$100\%$} \\
UpdraftNet ($n=8$)              & \rd{$16.4\%$}     & \rd{$75\,\text{s}$}  & $100\%$ \\
UpdraftNet ($n=16$)             & \rd{$\mathbf{16.1\%}$} & \rd{$37\,\text{s}$}  & $100\%$ \\
\rdisp Single soaring ($n=16$)    & \rd{$8.8\%$}           & \rd{$36\,\text{s}$}  & \rd{$100\%$} \\
\bottomrule
\end{tabular}
\end{table}

Fig.~\ref{fig:results}(f) reports saving as concurrent delivery tasks increase
at $n=8$. Saving grows from \rd{$3.5\%$ at 4 tasks to $22.6\%$ at 48 tasks} because more missions pass through confirmed
updraft corridors as the shared map densifies. Single-UAV
soaring \rd{grows more slowly ($0.8\%$ to $14.1\%$)} since its \rd{map is built from its own flights only}.

Table~\ref{tab:comparison} collects the full comparison at $n=8$ and
$n=16$. \rd{Map sharing raises the saving of the same eight UAVs from $9.2\%$ to $16.4\%$, whereas growing the swarm from 8 to 16 adds nothing ($16.1\%$), as the corridors each UAV uses are already mapped. Route voxels with at least two measurements in the UAV's own map: $30\%$ ($n=1$), $28\%$ ($n=8$, no sharing), $38\%$ ($n=8,16$, sharing).}
The observed \rd{16\%} savings fall below the theoretical upper
envelope for three reasons. Not all path voxels are updraft
cells, variance discounting suppresses poorly sampled regions,
and one-hop gossip propagation lags the aggregate swarm
observation rate.
\rd{These savings hold under the power model of Sec.~\ref{sec:framework}; at $n=16$ they fall to $4.4\%$ for $\beta=0.5$ and $0.6\%$ for $\beta=0.08\,\text{s/m}$. The model parameters are simulation assumptions needing hardware validation, and the wind is static within a trial.}

\textit{Simulation Limitations.}
To the best of our knowledge, no published dataset provides these values for delivery-class multirotors in building-induced urban updrafts.
We set them to physically reasonable values.
Hardware validation on a physical platform remains needed.

%% ------------------------------------------------------------------
\section{Conclusion}
In this work, we proposed a scheme, \textit{UpdraftNet}, a cooperative
framework for energy-aware urban UAV delivery that turns an existing
delivery fleet into a distributed wind-measurement network, where each
UAV shares local updraft readings via constant-size \rd{one-hop} \rd{messages} and routes missions over \rd{its growing locally fused} map. We established a coverage-time bound, a routing-optimality guarantee, a freshness-decay
characterisation, and a profitability condition for updraft detours.
Through simulations across 60 random urban configurations, UpdraftNet
achieved a \rd{$16.1\%$} reduction in cruise energy
at $n=16$ UAVs \rd{under the stated power model} versus shortest-path routing, with coverage
time dropping from \rd{$603\,\text{s}$} for a solo UAV to \rd{$37\,\text{s}$} for the full
fleet. In the future, we aim to explore adaptive gossip intervals
that respond to wind-field volatility, integrate solar-irradiance
maps to jointly exploit photovoltaic recharge and aerodynamic
lift, and co-optimise task assignments based on updraft-zone
coverage.

\section*{ACKNOWLEDGEMENT}
The authors gratefully acknowledge the financial support received from ANRF under the ARG--MATRICS scheme (File No. ANRF/ARGM/2025/002678/TS) for carrying out this research.

%% ------------------------------------------------------------------
\begin{thebibliography}{99}
\footnotesize

\bibitem{holmes}
J.\ D.\ Holmes, \emph{Wind Loading of Structures}, 3rd ed.
Boca Raton, FL, USA: CRC Press, 2018.

\bibitem{asce}
American Society of Civil Engineers, \emph{Minimum Design Loads and
Associated Criteria for Buildings and Other Structures}, ASCE/SEI 7-22.
Reston, VA, USA: ASCE, 2022.

\bibitem{allen}
M.\ J.\ Allen, ``Updraft model for development of autonomous soaring
uninhabited air vehicles,'' in \emph{Proc.\ 44th AIAA Aerosp.\ Sci.\ Meeting and
Exhibit}, Reno, NV, USA, Jan.\ 2006.

\bibitem{reichmann}
H.\ Reichmann, \emph{Cross-Country Soaring}, 7th ed. Santa Monica, CA, USA:
Thomson Publications/Soaring Society of America, 1993.

\bibitem{pooley}
A.\ Pooley, M.\ Gao, A.\ Sharma, S.\ Barnaby, Y.\ Gu, and J.\ Gross,
``Analysis of UAV thermal soaring via hawk-inspired swarm interaction,''
\emph{Biomimetics}, vol.~8, no.~1, art.~124, Mar.\ 2023.

\bibitem{wei}
X.\ Wei, Z.\ Yi, W.\ Li, L.\ Zhao, and W.\ Zhang, ``Energy harvesting fueling
the revival of self-powered unmanned aerial vehicles,'' \emph{Energy Convers.\
Manage.}, vol.~283, art.~116863, May 2023.

\bibitem{shakhatreh}
H.\ Shakhatreh, A.\ H.\ Sawalmeh, A.\ Al-Fuqaha, Z.\ Dou, E.\ Almaita,
I.\ Khalil, N.\ S.\ Othman, A.\ Khreishah, and M.\ Guizani, ``Unmanned
aerial vehicles (UAVs): A survey on civil applications and key research
challenges,'' \emph{IEEE Access}, vol.~7, pp.~48572--48634, 2019.

\bibitem{sarhan}
S.\ Sarhan, M.\ Rinaldi, S.\ Primatesta, and G.\ Guglieri, ``Noise-aware
UAV path planning in urban environment with reinforcement learning,''
\emph{Eng.\ Proc.}, vol.~90, no.~1, art.~3, Jan.\ 2025.

\bibitem{chen}
Y.\ Chen, R.\ Fu, and S.\ Jia, ``Flight path planning with ground risk and
noise mitigation for urban air mobility,'' \emph{Reliab.\ Eng.\ Syst.\ Saf.}, vol.~270,
art.~112176, 2026.

\bibitem{tan}
Q.\ Tan, Y.\ Li, H.\ Wu, P.\ Zhou, H.\ K.\ Lo, S.\ Zhong, and X.\ Zhang,
``Enhancing sustainable urban air transportation: Low-noise UAS flight
planning using noise assessment simulator,'' \emph{Aerosp.\ Sci.\ Technol.},
vol.~147, art.~109071, Apr.\ 2024.

\bibitem{aldao}
E.\ Aldao, L.\ M.\ Gonzalez-deSantos, H.\ Michinel, and H.\ Gonzalez-Jorge,
``UAV obstacle avoidance algorithm to navigate in dynamic
building environments,'' \emph{Drones}, vol.~6, no.~1, art.~16, Jan.\ 2022.

\bibitem{kawamura}
S.\ Kawamura, E.\ Kimoto, T.\ Fukushima, and Y.\ Taniike, ``Environmental
wind characteristics around the base of a tall building: A comparison
between model test and full scale experiment,'' \emph{J.\ Wind Eng.\ Ind.\
Aerodyn.}, vol.~28, pp.~149--158, 1988.

\bibitem{tominaga}
Y.\ Tominaga and M.\ Shirzadi, ``Wind tunnel measurement of three-dimensional turbulent flow structures around a building group: Impact
of high-rise buildings on pedestrian wind environment,'' \emph{Build.\ Environ.},
vol.~206, art.~108389, Dec.\ 2021.

\bibitem{britter}
R.\ E.\ Britter and S.\ R.\ Hanna, ``Flow and dispersion in urban areas,''
\emph{Annu.\ Rev.\ Fluid Mech.}, vol.~35, pp.~469--496, Jan.\ 2003.

\bibitem{la}
H.\ M.\ La and W.\ Sheng, ``Distributed sensor fusion for scalar field
mapping using mobile sensor networks,'' \emph{IEEE Trans.\ Cybern.}, vol.~43,
no.~2, pp.~766--778, Apr.\ 2013.

\bibitem{bayerlein}
H.\ Bayerlein, M.\ Theile, M.\ Caccamo, and D.\ Gesbert, ``Multi-UAV path
planning for wireless data harvesting with deep reinforcement learning,''
\emph{IEEE Open J.\ Commun.\ Soc.}, vol.~2, pp.~1171--1187, Jun.\ 2021.

\bibitem{boyd}
S.\ Boyd, A.\ Ghosh, B.\ Prabhakar, and D.\ Shah, ``Randomized gossip
algorithms,'' \emph{IEEE Trans.\ Inf.\ Theory}, vol.~52, no.~6, pp.~2508--2530,
Jun.\ 2006.

\bibitem{motwani}
R.\ Motwani and P.\ Raghavan, \emph{Randomized Algorithms}. Cambridge,
U.K.: Cambridge Univ.\ Press, 1995.

\bibitem{doerr}
B.\ Doerr, ``Probabilistic tools for the analysis of randomized optimization
heuristics,'' in \emph{Theory of Evolutionary Computation}, B.\ Doerr and F.\
Neumann, Eds. Cham, Switzerland: Springer, 2020, ch.~1.

\bibitem{thrun}
S.\ Thrun, W.\ Burgard, and D.\ Fox, \emph{Probabilistic Robotics}. Cambridge,
MA, USA: MIT Press, 2005.

\bibitem{lawrance}
N.\ R.\ J.\ Lawrance and S.\ Sukkarieh, ``Autonomous exploration of a
wind field with a gliding aircraft,'' \emph{J.\ Guid.\ Control Dyn.}, vol.~34, no.~3,
pp.~719--733, May 2011.

\bibitem{hart}
P.\ E.\ Hart, N.\ J.\ Nilsson, and B.\ Raphael, ``A formal basis for the heuristic
determination of minimum cost paths,'' \emph{IEEE Trans.\ Syst.\ Sci.\ Cybern.},
vol.~4, no.~2, pp.~100--107, Jul.\ 1968.

\bibitem{kirschstein}
T.\ Kirschstein, ``Comparison of energy demands of drone-based and
ground-based parcel delivery services,'' \emph{Transp.\ Res.\ Part D: Transp.\
Environ.}, vol.~78, art.~102209, 2020.

\bibitem{leishman}
\rd{J.\ G.\ Leishman, \emph{Principles of Helicopter Aerodynamics}, 2nd ed.
Cambridge, U.K.: Cambridge Univ.\ Press, 2006.}

\bibitem{newman}
\rd{D.\ J.\ Newman and L.\ Shepp, ``The double dixie cup problem,''
\emph{Amer.\ Math.\ Monthly}, vol.~67, no.~1, pp.~58--61, Jan.\ 1960.}

\end{thebibliography}
\end{document}